\section*{Computational verification and reproducibility}

The computational proof consists of explicit finite rational identities.
The supplementary material provides the seven-vertex graph atlas, the
rooted flag spaces, primitive pattern counts, three rational witness
files, and their generators and verifiers. The main certificate contains
$184$ positive square terms, $37$ unrestricted identity columns, and
$826$ positive host remainders. The multipliers of $\Phi_\triangle$ and
$J_{\mathrm{del}}$ are positive, while those of $J_0$ and $J_1$ are zero.
Every one of the $1,044$ host identities is checked with exact integer or
rational arithmetic. A second implementation contracts integer Gram
matrices against independently recounted host tensors. The equality-span
witnesses are bound to the same main certificate and checked coordinate
by coordinate. Numerical optimization is used only to discover witnesses;
its floating-point output is not accepted as proof.

The supplementary programs also enumerate all $65,536$ Boolean Clebsch
profiles and all $6,496$ small-offset presentations under the convention
in Section~14.1. An independent enumeration of labeled exceptional
vertices gives the same canonical family. Formula (75) is recomputed for
every presentation, and the distinct polynomial formulas are additionally
checked on actual finite graphs. The branch counts in Section~14.1 have
a separate direct census check. These are verifications of the finite
inputs and construction formulas. The global bound, stability arguments,
and passage to eventual extremality are proved in the text. In
particular, the computations do not provide a numerical global threshold
$k_0(j)$.

The command \texttt{python3 src/verify\_project.py} reruns the complete
finite verification chain. Definitions, normalization conventions,
dependencies, and individual commands are specified in the accompanying
\texttt{README.md} and computational reproduction guide.

\paragraph{AI-assisted research methodology.}
ChatGPT (OpenAI) assisted with proof exploration, symbolic calculations,
implementation of generators and verifiers, and internal mathematical
cross-checks. Its mathematical output was treated as candidate reasoning,
not as a certificate of correctness. The supplied deterministic programs
perform the exact finite checks described above, and the analytic
arguments are given explicitly in the paper. 